\section{Application I: Manifesto Project RILE Scoring}\label{sec:manifesto}

Our first empirical application demonstrates OPS on party manifestos from the Manifesto Project, estimating parties' positions on the left-right (RILE) dimension. This task is central to comparative politics and provides a well-defined oracle with established benchmarks.

\subsection{Task Definition}

The Manifesto Project \citep{manifesto2021} codes party positions by classifying sentences into policy categories. The RILE scale aggregates these codes into a single left-right score:
\[
\text{RILE} = \frac{R - L}{R + L + N} \times 100,
\]
where $R$ is the count of right-leaning statements, $L$ is left-leaning, and $N$ is neutral. Scores range from $-100$ (far left) to $+100$ (far right).

\paragraph{Oracle Definition.}
We define the oracle $\fstar$ as the RILE score that would be assigned by expert coders with access to the full manifesto. In practice, we approximate $\fstar$ using:
\begin{enumerate}[nosep]
    \item Published Manifesto Project expert scores (when available)
    \item High-quality deliberation model (GPT-4o with detailed rubric)
    \item Human validation on a sample
\end{enumerate}

\paragraph{Rubric Design.}
The summarization rubric specifies what information must be preserved for RILE scoring. Following the Manifesto Project codebook, the rubric includes:
\begin{itemize}[nosep]
    \item Policy positions on economic issues (welfare, taxation, regulation)
    \item Policy positions on social issues (immigration, law and order, traditional values)
    \item Positions on political system (decentralization, democracy, EU integration)
    \item Party's self-positioning relative to competitors
\end{itemize}

The \texttt{ThinkingTrees} implementation is in \texttt{src/tasks/manifesto/rubrics.py}.

\subsection{Experimental Setup}

\paragraph{Dataset.}
We use manifestos from the Manifesto Project corpus spanning:
\begin{itemize}[nosep]
    \item Countries: [TBD: specify countries]
    \item Years: [TBD: specify range]
    \item Parties: [TBD: number of parties]
    \item Total documents: [TBD]
    \item Average length: [TBD tokens]
\end{itemize}

\paragraph{Baselines.}
We compare three approaches:
\begin{enumerate}[nosep]
    \item \textbf{Direct scoring}: Apply the oracle directly to full manifestos (gold standard but expensive)
    \item \textbf{Truncation}: Truncate to context window and score
    \item \textbf{OPS}: Hierarchical summarization with audit
\end{enumerate}

\paragraph{Metrics.}
\begin{itemize}[nosep]
    \item Correlation with expert scores (Pearson's $r$)
    \item Mean absolute error (MAE) on RILE scale
    \item Audit violation rates ($\hat{p}_{\text{suff}}, \hat{p}_{\text{assoc}}, \hat{p}_{\text{idem}}$)
    \item Compression ratio (original tokens / summary tokens)
\end{itemize}

\subsection{Results}

[TBD: Present results including:]

\paragraph{Audit Statistics.}
Table showing empirical violation rates for each condition, with confidence intervals.

\paragraph{Correlation with Expert Scores.}
Comparison of OPS vs. baselines on correlation with published Manifesto Project scores.

\paragraph{Compression vs. Preservation Trade-off.}
Plot showing how correlation changes as a function of compression ratio.

\paragraph{Error Analysis.}
Examples of documents where OPS fails and why (e.g., complex cross-references, implicit positions).

\subsection{Comparison with Benoit et al. (2025)}

\citet{BenoitEtAl2025} use LLM ensembles to estimate party positions, achieving high correlation with expert surveys. Our approach differs in several ways:

\begin{itemize}[nosep]
    \item \textbf{Theoretical guarantees}: OPS provides provable bounds on information loss
    \item \textbf{Auditing}: We measure and report violation rates, enabling quality control
    \item \textbf{Scalability}: Hierarchical summarization handles arbitrary document lengths
    \item \textbf{Preference learning}: Our framework extends to training policies on summaries
\end{itemize}

[TBD: Empirical comparison if data overlap permits]

\subsection{Discussion}

\paragraph{When OPS Helps.}
OPS is most beneficial when:
\begin{itemize}[nosep]
    \item Documents exceed context windows (manifestos often do)
    \item Multiple downstream analyses use the same summaries
    \item Auditing and quality control are important
    \item Preferences need to be collected from compressed representations
\end{itemize}

\paragraph{Rubric Design Lessons.}
[TBD: Lessons learned from rubric iteration]

\paragraph{Limitations.}
[TBD: Task-specific limitations]
